\SourceLeaf{186}
\SeventhArticle{11}
Whenever some or all of $a, a', a''$, etc., are irrational, it will be useful to transform the rational function from which we derive the numbers $R, R', R''$, etc., into a polynomial. Since the principles of such a transformation are not found in algebraic textbooks, we shall explain them briefly here. Given three polynomials $Z,\zeta,\zeta'$ in the indeterminate $z$, we seek a polynomial that can take the place of the fraction $\frac Z\zeta$ when $z$ is taken as any root of the equation $\zeta'=0$. We shall assume, however, that $\zeta$ vanishes for none of these values of $z$, or equivalently, that $\zeta$ and $\zeta'$ have no common nonconstant divisor. We shall denote the exponents of the highest powers of $z$ in $\zeta$ and $\zeta'$ by $k,k'$.

Divide $\zeta$ by $\zeta'$ in the usual way until the degree of the remainder is less than $k'$. Let the remainder be $=\frac{\zeta''}{\lambda}$ and its degree $=k''$, so that $\frac1\lambda z^{k''}$ is its highest term; we shall put the quotient of the division $=\frac p\lambda$. Likewise, let the division of the function $\zeta'$ by $\zeta''$ give the remainder $\frac{\zeta'''}{\lambda'}$ of degree $k'''$ and the quotient $\frac{p'}{\lambda'}$; then let the division of the function $\zeta''$ by $\zeta'''$ give the remainder $\frac{\zeta''''}{\lambda''}$ of degree $k''''$ and the quotient $\frac{p''}{\lambda''}$, and so on, until in the sequence of functions $\zeta'',\ \zeta'''.\ \zeta''''$, etc., each of which has $1$ as the coefficient of its highest term, we reach $\zeta^{(m)}=1$. It is easy to see that this must eventually happen, since none of the functions $\zeta,\ \zeta'\ \zeta'',\ \zeta'''$, etc., can have a common nonconstant divisor with the preceding one. Thus division without a remainder certainly cannot occur while the degree of the divisor is greater than $0$. We shall therefore have the sequence of equations
\[
\begin{aligned}
\zeta''&=\lambda\zeta-p\zeta'\\
\zeta'''&=\lambda'\zeta'-p'\zeta''\\
\zeta''''&=\lambda''\zeta''-p''\zeta'''\\
\zeta'''''&=\lambda'''\zeta'''-p'''\zeta''''\\
&\text{etc., up to}\\
\zeta^{(m)}&=\lambda^{(m-2)}\zeta^{(m-2)}-p^{(m-2)}\zeta^{(m-1)}
\end{aligned}
\]
where $\zeta'',\allowbreak \ \zeta''',\allowbreak \ \zeta''''\ldotp\ldotp\ldotp\ldotp\zeta^{(m)}$ are polynomials in $z$ of degrees $k'',\allowbreak \ k''',\allowbreak \ k''''\allowbreak\ldotp\ldotp\ldotp\ldotp k^{(m)}$; the numbers $k',\allowbreak \ k'',\allowbreak \ k'''\ldotp\ldotp\ldotp\ldotp k^{(m)}$ decrease steadily down to the last, $k^{(m)}=0$; and $p,\allowbreak p',\allowbreak p'',\allowbreak p'''$, etc., are also polynomials in $z$, of degrees $k-k',\allowbreak \ k'-k'',\allowbreak \ k''-k''',\allowbreak \ k'''-k''''$, etc. (except when $k<k'$, in which case $p=0$ must evidently be set).

With these preparations, we form a \textit{second} sequence of polynomials in $z$, namely $\eta,\ \eta',\ \eta'',\ \eta'''\ldotp\ldotp\ldotp\ldotp\eta^{(m)}$. We shall set $\eta=1,\ \eta'=0$, and derive each of the remaining terms\SourceLeaf{187} from the two preceding ones by the same rule that connects the functions $\zeta,\zeta',\zeta'',\zeta'''$, etc., namely by the equations
\[
\begin{aligned}
\eta''&=\lambda\eta-p\eta'\\
\eta'''&=\lambda'\eta'-p'\eta''\\
\eta''''&=\lambda''\eta''-p''\eta'''\\
\eta'''''&=\lambda'''\eta'''-p'''\eta''''\quad\text{etc., up to}\\
\eta^{(m)}&=\lambda^{(m-2)}\eta^{(m-2)}-p^{(m-2)}\eta^{(m-1)}
\end{aligned}
\]
Evidently $\eta''=\lambda$ here has degree $0$; $\eta'''=-\lambda p'$ has degree $k'-k''$; and likewise the following functions $\eta'''',\eta'''''$, etc., have degrees $k'-k''',\ k'-k''''$, etc., respectively, so that the last, $\eta^{(m)}$, reaches degree $k'-k^{(m-1)}$.

Let us further consider the \textit{third} sequence of functions, $\zeta-\zeta\eta,\ \zeta'-\zeta\eta',\ \zeta''-\zeta\eta'',\ \zeta'''-\zeta\eta'''$, etc. A similar relation evidently holds between any three successive terms of this sequence, namely
\[
\begin{aligned}
\zeta''-\zeta\eta''&=\lambda(\zeta-\zeta\eta)-p(\zeta'-\zeta\eta')\\
\zeta'''-\zeta\eta'''&=\lambda'(\zeta'-\zeta\eta')-p'(\zeta''-\zeta\eta'')\\
\zeta''''-\zeta\eta''''&=\lambda''(\zeta''-\zeta\eta'')-p''(\zeta'''-\zeta\eta''')
\end{aligned}
\]
The first of these functions gives $=0$, the second $=\zeta'$; hence it is easily inferred that each is divisible by $\zeta'$.

It now follows without difficulty that the fraction $\frac Z\zeta$ can be replaced by the polynomial $Z\eta^{(m)}$, provided that $z$ is assigned only values that are roots of the equation $\zeta'=0$. Indeed, the difference $\frac{Z(1-\zeta\eta^{(m)})}{\zeta}$ must vanish for such a value of $z$, since $1-\zeta\eta^{(m)}=\zeta^{(m)}-\zeta\eta^{(m)}$ is divisible by $\zeta'$.

In place of the function $Z\eta^{(m)}$, one may also take its remainder on division by $\zeta'$; the degree of this remainder will be less than the degree of the \mbox{function $\zeta'$}.

This remainder, however, can be obtained directly and more conveniently by the following algorithm. Form the equations
\[
\begin{aligned}
Z&=q'\zeta'+Z'\\
Z'&=q''\zeta''+Z''\\
Z''&=q'''\zeta'''+Z'''\\
Z'''&=q''''\zeta''''+Z''''\quad\text{etc., up to}\\
Z^{(m-1)}&=q^{(m)}\zeta^{(m)}+Z^{(m)}
\end{aligned}
\]
\SourceLeaf{188}by dividing $Z$ successively by $\zeta'$, then the remainder of the first division, $Z'$, by $\zeta''$, then the remainder of the second division by $\zeta'''$, and so on. Since the remainder always has a lower degree than the divisor, the degrees of the functions $Z',Z'',Z''',Z''''$, etc., will be less than $k',k'',k''',k''''$, etc., respectively. The last, $Z^{(m)}$, necessarily becomes $=0$, since the divisor $\zeta^{(m)}$ $=1$. We therefore have
\[
Z=q'\zeta'+q''\zeta''+q'''\zeta'''+q''''\zeta''''+\text{etc.}+q^{(m)}\zeta^{(m)}
\]
When $z$ is taken only as a root of the equation $\zeta'=0$, we have $\zeta'=0$, $\zeta''=\zeta\eta''$, $\zeta'''=\zeta\eta'''$, $\zeta''''=\zeta\eta''''$, etc.; hence, under the same restriction,
\[
\frac{Z}{\zeta}=q''\eta''+q'''\eta'''+q''''\eta''''+\text{etc.}+q^{(m)}\eta^{(m)}
\]
The degree of this expression must be less than $k'$. For, since the degrees of the quotients $q'',q''',q''''$, etc., must be less than $k'-k''$, $k''-k'''$, $k'''-k''''$, etc., the degrees of the individual parts $q''\eta''$, $q'''\eta'''$, $q''''\eta''''$, etc., will be less than $k'-k''$, $k'-k'''$, $k'-k''''$, etc.

Finally, we observe that if, among the values of the indeterminate $z$ to be substituted in the fraction $\frac{Z}{\zeta}$, rational values happen to occur together with irrational ones, it will be more advantageous to separate the former from the latter and include only the latter in the equation $\zeta'=0$. For rational values, such a shortcut in the calculation will not be needed; for irrational values, however, the calculation will be simpler, the lower the degree of the polynomial to which the rational function can be reduced.

\SeventhArticle{12}
Here is an example of the transformation explained in the preceding article. Let the rational function be
\[
\frac{z^6-\frac{50}{39}z^4+\frac{283}{715}zz-\frac{256}{15015}}{7z^6-\frac{105}{13}z^4+\frac{315}{143}zz-\frac{35}{429}}
\]
where $z$ stands for any root of the equation
\[
z^7-\frac{21}{13}z^5+\frac{105}{143}z^3-\frac{35}{429}z=0
\]
If we wished to include all seven roots here, we would arrive at a polynomial of degree six. But for the rational value $z=0$, the calculation of the fraction is evidently immediate and gives the value $\frac{256}{1225}$. Thus, setting this root aside, we shall confine ourselves to the equation of degree six:\SourceLeaf{189}
\[
z^6-\frac{21}{13}z^4+\frac{105}{143}zz-\frac{35}{429}=0
\]
In this way, we can readily foresee that a polynomial of degree four will result. Applying the preceding rules now gives the following:
\[
\begin{aligned}
\zeta&=7z^6-\frac{105}{13}z^4+\frac{315}{143}zz-\frac{35}{429}\\[5pt]
\zeta'&=z^6-\frac{21}{13}z^4+\frac{105}{143}zz-\frac{35}{429}\\[5pt]
\zeta''&=z^4-\frac{10}{11}zz+\frac{5}{33}\\[5pt]
\zeta'''&=zz-\frac{3}{7}\\[5pt]
\zeta''''&=1
\end{aligned}\]
\[\begin{aligned}
\lambda&=\frac{13}{42}&p&=\frac{13}{6}\\[5pt]
\lambda'&=-\frac{4719}{280}&p'&=-\frac{4719}{280}zz+\frac{3333}{280}\\[5pt]
\lambda''&=-\frac{147}{8}&p''&=-\frac{147}{8}zz+\frac{777}{88}
\end{aligned}\]
\[\begin{aligned}
\eta&=1\\[5pt]
\eta'&=0\\[5pt]
\eta''&=\frac{13}{42}\\[5pt]
\eta'''&=\frac{20449}{3920}zz-\frac{14443}{3920}\\[5pt]
\eta''''&=\frac{61347}{640}z^4-\frac{127413}{1120}zz+\frac{120263}{4480}
\end{aligned}\]
\[\begin{aligned}
Z&=z^6-\frac{50}{39}z^4+\frac{283}{715}zz-\frac{256}{15015};&q'&=1\\[5pt]
Z'&=\frac{1}{3}z^4-\frac{22}{65}zz+\frac{323}{5005}&q''&=\frac{1}{3}\\[5pt]
Z''&=-\frac{76}{2145}zz+\frac{632}{45045}&q'''&=-\frac{76}{2145}\\[5pt]
Z'''&=-\frac{4}{3465}&q''''&=-\frac{4}{3465}
\end{aligned}
\]
Finally, the following polynomial, equivalent to the given fraction, is obtained:
\[
-\frac{1859}{16800}z^4-\frac{1573}{29400}zz+\frac{7947}{39200}
\]

